Fugitive emissions arise across segments national production statistics
alone cannot represent. We develop an infrastructure-constrained inventory for Pakistan from
the 2006 IPCC Guidelines and the segment-resolved 2019 Refinement, adding every segment
public data supports. Gas balances, Pakistan's Common Reporting Table
(CRT), flare records and methane-plume detections provide separate evaluation layers.

On common FY2023--24 production support, legacy and refined methods give
\SI{3.317}{} and \SI{4.485}{\mega\tonne} \COe{}, respectively. The complete public-data
ledger gives an \emph{upper-practice conservative (default-practice) configuration} of
\SI{23.499}{\mega\tonne}---not an empirically centred best estimate: segment
coverage matters more than the factor revision. Joint structural configurations
span \SIrange{9.601}{24.694}{\mega\tonne}; the reference sits high in that
envelope---no probability attaches to any point within---and the joint floor
ships as \emph{fugitives\_joint\_minimum}. Uncertainty is dominated by
gas-distribution practice and a transferred post-meter industrial/power factor.

Adjacent-fiscal-year reconstructions bracket calendar 2022 at
\SIrange{23.761}{25.187}{\mega\tonne} against
\SI{4.749}{\mega\tonne} in Pakistan's 2022 CRT on a common GWP basis. The official arithmetic
reproduces exactly; the difference is activity, factor and boundary choices. Across
the nine-year panel, declining upstream emissions are partly offset by gas-network
growth. Matched detections occur near fields with \SI{71.1}{\percent} of reported supply,
yet production only weakly rank-associates with observed activity.

The framework identifies two measurement priorities---distribution technology and post-meter
leakage---and reconciles energy statistics, inventory reporting and independent
observations in data-scarce systems.
